\section{第7章の解答}\label{解答:第7章}

\paragraph{7.1}
$x$ は微分すると定数になり,\,$\e^{3x}$ は積分しても指数関数のままである.したがって
\[
 u=x,\qquad v'=\e^{3x},\qquad
 u'=1,\qquad v=\frac13\e^{3x}
\]
と選ぶ.部分積分により
\[
\begin{aligned}
 \int x\e^{3x}\dd x
 &=\frac x3\e^{3x}-\frac13\int\e^{3x}\dd x\\
 &=\frac x3\e^{3x}-\frac19\e^{3x}+C\\
 &=\frac{3x-1}{9}\e^{3x}+C.
\end{aligned}
\]
$\e^{3x}$ を積分するたびに $3$ で割るため,最後の項の係数は $\frac19$ になる.右辺を微分すれば
\[
\begin{aligned}
 \frac{\dd}{\dd x}\left(\frac{3x-1}{9}\e^{3x}\right)
 &=\frac13\e^{3x}+\frac{3x-1}{9}\cdot3\e^{3x}\\
 &=\frac{1+3x-1}{3}\e^{3x}\\
 &=x\e^{3x}
\end{aligned}
\]
となる.

\paragraph{7.2}
$\log x$ を微分すれば $\frac1x$ になり,\,$x^2$ の原始関数 $\frac{x^3}{3}$ と掛けると整式が残る.そこで $u=\log x$,\,$v'=x^2$ と選ぶ.
\[
\begin{aligned}
 \int x^2\log x\dd x
 &=\frac{x^3}{3}\log x
  -\int\frac{x^3}{3}\cdot\frac1x\dd x\\
 &=\frac{x^3}{3}\log x-\frac13\int x^2\dd x\\
 &=\frac{x^3}{3}\log x-\frac{x^3}{9}+C.
\end{aligned}
\]
$x^2$ を微分する側に選ぶこともできるが,その場合は $\log x$ の原始関数を求めたうえで,残る積にも対数が現れる.上の選択では1回の部分積分で対数が積分記号の中からなくなるため,見通しがよい.答えを微分すると,\,$\frac{x^3}{3}\cdot\frac1x=\frac{x^2}{3}$ と $-\frac{x^3}{9}$ の微分が打ち消し合う.

\paragraph{7.3}
整式を微分する側に選ぶと,最初の部分積分で三次式が二次式になる.
\[
 \int x^3\sin x\dd x
 =-x^3\cos x+3\int x^2\cos x\dd x.
\]
続いて同じ方針で
\[
 \int x^2\cos x\dd x
 =x^2\sin x-2\int x\sin x\dd x
\]
とし,二次式を一次式にする.さらに
\[
\begin{aligned}
 \int x\sin x\dd x
 &=-x\cos x+\int\cos x\dd x\\
 &=-x\cos x+\sin x+C
\end{aligned}
\]
とすれば,整式の因子は定数になった.これらを逆の順序に戻すと
\[
\begin{aligned}
 \int x^3\sin x\dd x
 &=-x^3\cos x
   +3\{x^2\sin x-2(-x\cos x+\sin x)\}+C\\
 &=-x^3\cos x+3x^2\sin x+6x\cos x-6\sin x+C.
\end{aligned}
\]
微分すると,\,$x^3\sin x$ 以外の項は
\[
 -3x^2\cos x+3x^2\cos x,
 \quad 6x\sin x-6x\sin x,
 \quad 6\cos x-6\cos x
\]
の組になって消える.次数は三次,二次,一次,定数と1つずつ下がっている.

\paragraph{7.4}
本文と同じく,積分定数を固定するため
\[
 I(x)=\int_0^x\e^{2t}\cos t\dd t,
 \qquad
 J(x)=\int_0^x\e^{2t}\sin t\dd t
\]
と定める.どちらも指数関数を積分する側に選ぶと,
\[
\begin{aligned}
 I(x)
 &=\left[\frac12\e^{2t}\cos t\right]_0^x
   +\frac12\int_0^x\e^{2t}\sin t\dd t\\
 &=\frac12\e^{2x}\cos x-\frac12+\frac12J(x),\\
 J(x)
 &=\left[\frac12\e^{2t}\sin t\right]_0^x
   -\frac12\int_0^x\e^{2t}\cos t\dd t\\
 &=\frac12\e^{2x}\sin x-\frac12I(x).
\end{aligned}
\]
第二式を第一式へ代入して
\[
 I(x)=\frac12\e^{2x}\cos x-\frac12
       +\frac14\e^{2x}\sin x-\frac14I(x).
\]
したがって
\[
 \frac54I(x)=\frac14\e^{2x}(2\cos x+\sin x)-\frac12
\]
となり,両辺に $\frac45$ を掛けて
\[
 I(x)=\frac15\e^{2x}(2\cos x+\sin x)-\frac25
\]
を得る.定数をまとめれば
\[
 \boxed{\int\e^{2x}\cos x\dd x
 =\frac15\e^{2x}(2\cos x+\sin x)+C}
\]
となる.右辺を微分すると
\[
\begin{aligned}
 &\frac15\e^{2x}
 \{2(2\cos x+\sin x)-2\sin x+\cos x\}\\
 &\qquad=\frac15\e^{2x}\cdot5\cos x\\
 &\qquad=\e^{2x}\cos x
\end{aligned}
\]
である.\,$\e^{2x}$ の原始関数に付く $\frac12$ を,2回とも保つことが必要である.

\paragraph{7.5}
$t=\log x$ と置く.\,$x>0$ だから $x=\e^t$ と戻せて,\,$\dd x=\e^t\dd t$ となる.したがって
\[
 \int\cos(\log x)\dd x
 =\int\e^t\cos t\dd t.
\]
右辺では $\cos t$ を微分し,\,$\e^t$ を積分する.本文で得た $\int\e^t\sin t\dd t$ の結果も使うと
\[
\begin{aligned}
 \int\e^t\cos t\dd t
 &=\e^t\cos t+\int\e^t\sin t\dd t\\
 &=\e^t\cos t+\frac12\e^t(\sin t-\cos t)+C\\
 &=\frac12\e^t(\sin t+\cos t)+C.
\end{aligned}
\]
よって元の変数へ戻して
\[
 \int\cos(\log x)\dd x
 =\frac x2\{\sin(\log x)+\cos(\log x)\}+C.
\]
微分すると,括弧をそのまま残す項と,括弧を微分する項が加わる.前者は $\frac12\{\sin(\log x)+\cos(\log x)\}$,後者は $\frac12\{\cos(\log x)-\sin(\log x)\}$ なので,和は $\cos(\log x)$ になる.

\paragraph{7.6}
$t=1+x^2$ と置くと,\,$\dd t=2x\dd x$ であり,端点は $x=0,1$ から $t=1,2$ へ移る.したがって
\[
 \int_0^1x\log(1+x^2)\dd x
 =\frac12\int_1^2\log t\dd t.
\]
ここで $\log t$ を微分し,\,$1$ を積分する部分積分を使う.
\[
\begin{aligned}
 \frac12\int_1^2\log t\dd t
 &=\frac12\left\{[t\log t]_1^2-\int_1^2 1\dd t\right\}\\
 &=\frac12(2\log2-1)\\
 &=\log2-\frac12.
\end{aligned}
\]
もとの式にある $x$ は,対数を微分するために外に残したのではなく,置換の $x\dd x=\frac12\dd t$ として使った.置換後は対数1つの積分になるので,そこで部分積分を行えばよい.

\paragraph{7.7}
$I=\int_0^1\sqrt{4+x^2}\dd x$ と置く.根号を微分し,\,$1$ を積分すると,
\[
 I=[x\sqrt{4+x^2}]_0^1
   -\int_0^1\frac{x^2}{\sqrt{4+x^2}}\dd x.
\]
分子を $x^2=(4+x^2)-4$ と分けると
\[
 \frac{x^2}{\sqrt{4+x^2}}
 =\sqrt{4+x^2}-\frac4{\sqrt{4+x^2}}.
\]
よって,\,$L=\int_0^1\frac1{\sqrt{4+x^2}}\dd x$ と書けば
\[
 I=\sqrt5-I+4L
\]
となり,\,$I$ を移項して
\[
 2I=\sqrt5+4L
\]
を得る.
残る $L$ では,\,$r(x)=\sqrt{4+x^2}$ とすると $r'(x)=\frac{x}{r(x)}$ である.また $r(x)>|x|$ なので $x+r(x)>0$ が成り立ち,
\[
\begin{aligned}
 \frac{\dd}{\dd x}\log\{x+r(x)\}
 &=\cfrac{1+\cfrac{x}{r(x)}}{x+r(x)}\\
 &=\frac1{r(x)}
\end{aligned}
\]
となる.したがって
\[
\begin{aligned}
 L&=[\log(x+\sqrt{4+x^2})]_0^1\\
  &=\log(1+\sqrt5)-\log2\\
  &=\log\frac{1+\sqrt5}{2}.
\end{aligned}
\]
これを戻して
\[
 \boxed{\int_0^1\sqrt{4+x^2}\dd x
 =\frac{\sqrt5}{2}+2\log\frac{1+\sqrt5}{2}}
\]
を得る.下端 $x=0$ での対数の値は $\log2$ であり,\,$0$ にはならない.被積分関数が $\sqrt{1+x^2}$ なら下端での値は $\log1=0$ となるので,その場合との違いに注意する.

\paragraph{7.8}
$u=x$,\,$v'=f'(x)$ とすれば $u'=1$,\,$v=f(x)$ を使える.したがって
\[
\begin{aligned}
 \int_0^1xf'(x)\dd x
 &=[xf(x)]_0^1-\int_0^1f(x)\dd x\\
 &=2-\frac32\\
 &=\frac12.
\end{aligned}
\]
もう1つは $u=1-x$,\,$v'=f'(x)$ とする.今度は $u'=-1$ だから
\[
\begin{aligned}
 \int_0^1(1-x)f'(x)\dd x
 &=[(1-x)f(x)]_0^1+\int_0^1f(x)\dd x\\
 &=(0-1)+\frac32\\
 &=\frac12.
\end{aligned}
\]
端点の積と $f$ の定積分が分かっているため,\,$f$ の具体的な式は不要であった.2つの結果を加えると $1$ になり,
\[
\begin{aligned}
 \int_0^1\{x+(1-x)\}f'(x)\dd x
 &=\int_0^1f'(x)\dd x\\
 &=f(1)-f(0)\\
 &=1
\end{aligned}
\]
とも一致する.
